\chapter{函数的收敛}

\begin{solution}{essential-uniform-equivalence}
若存在零测集 \(N\) 使
\(\sup_{X\setminus N}|f_n-f|\to0\)，则本质上确界不超过该上确界。
反过来，对每个 \(n,k\)，由本质上确界定义可删去零测集 \(N_{n,k}\)，使
\[
 |f_n-f|\le\operatorname*{ess\,sup}|f_n-f|+1/k
\]
在其补集成立。令 \(N=\cup_{n,k}N_{n,k}\)，仍为零测集。对每个 \(n\)
令 \(k\to\infty\)，得到
\[
 \sup_{X\setminus N}|f_n-f|
 \le\operatorname*{ess\,sup}|f_n-f|,
\]
故一致收敛。
\end{solution}

\begin{solution}{measure-metric-triangle}
若 \(d_\mu(f,g)<r\)、\(d_\mu(g,h)<s\)，可取
\(r'<r,s'<s\) 略大于相应下确界，使
\[
 \mu(|f-g|>r')\le r',\qquad
 \mu(|g-h|>s')\le s'.
\]
若 \(r'+s'>1\)，则 \(d_\mu(f,h)\le1<r'+s'\)，无需再证。以下设
\(r'+s'\le1\)。由
\(\{|f-h|>r'+s'\}\subset\{|f-g|>r'\}\cup\{|g-h|>s'\}\)，
得到
\[
 d_\mu(f,h)\le r'+s'.
\]
令 \(r'\downarrow d_\mu(f,g)\)、\(s'\downarrow d_\mu(g,h)\)。
若 \(d_\mu(f,g)=0\)，则对每个 \(k\) 可得
\(\mu(|f-g|>1/k)=0\)，所以 \(f=g\) 几乎处处；反向显然。
\end{solution}

\begin{solution}{measure-convergence-cauchy}
依测度 Cauchy 条件是：对任意 \(\varepsilon,\eta>0\)，存在 \(N\)，使
\[
 m,n\ge N\Longrightarrow
 \mu(\{|f_m-f_n|>\varepsilon\})<\eta.
\]
若 \(d_\mu(f_m,f_n)\to0\)，取距离小于
\(\min\{\varepsilon,\eta\}/2\) 即得该条件。反之，取
\(\varepsilon=\eta=\delta\)，最终
\(\mu(|f_m-f_n|>\delta)\le\delta\)，所以
\(d_\mu(f_m,f_n)\le\delta\)。
\end{solution}

\begin{solution}{measure-limit-operations}
若 \(\Phi\) 一致连续，给定 \(\varepsilon>0\)，取 \(\delta>0\) 使
\(|u-v|\le\delta\) 蕴含
\(|\Phi(u)-\Phi(v)|\le\varepsilon\)。于是
\[
 \{|\Phi(f_n)-\Phi(f)|>\varepsilon\}
 \subseteq\{|f_n-f|>\delta\},
\]
右侧测度趋于零。

以下设 \(\mu(X)<\infty\)。因为 \(f,g\) 几乎处处有限，可先选 \(M\) 使
\[
 \mu(\{|f|>M\}\cup\{|g|>M\})
\]
任意小。连续函数在 \([-M-1,M+1]\) 上一致连续，而依测度收敛使
\(\{|f_n-f|>1\}\) 的测度趋零，故在其余部分可以使用紧区间上的一致连续
估计。这证明一般连续复合。

对乘法，写
\[
 f_ng_n-fg=f_n(g_n-g)+g(f_n-f).
\]
在上述小例外集之外，最终有 \(|f_n|\le M+1\)、\(|g|\le M\)。再除去
\[
 \left\{|g_n-g|>\frac{\varepsilon}{2(M+1)}\right\}
 \quad\text{和}\quad
 \left\{|f_n-f|>\frac{\varepsilon}{2\max\{M,1\}}\right\},
\]
乘积差便不超过 \(\varepsilon\)。各坏集的测度都可任意小，故
\(f_ng_n\to fg\) 依测度。
\end{solution}

\begin{solution}{ae-to-measure-finite}
固定 \(\varepsilon>0\)，令
\[
 E_N=\bigcup_{n\ge N}\{|f_n-f|>\varepsilon\}.
\]
则 \(E_N\downarrow E\)，而 a.e. 收敛说明 \(E\) 零测。因
\(\mu(E_1)\le\mu(X)<\infty\)，测度关于递减集合列连续，故
\(\mu(E_N)\downarrow0\)。又
\(\{|f_N-f|>\varepsilon\}\subset E_N\)，得到依测度收敛。
\end{solution}

\begin{solution}{typewriter-details}
令 \(n=2^m+k\)，其中 \(m\ge0,0\le k<2^m\)，并取
\[
 f_n=\mathbf1_{[k2^{-m},(k+1)2^{-m})}.
\]
则 \(\|f_n\|_p^p=2^{-m}\to0\)。对非二进端点 \(x\)，每个层级 \(m\)
恰有一个 \(k\) 使 \(f_{2^m+k}(x)=1\)，同一层另有许多项在 \(x\) 取
零，所以函数值既无限次为 \(1\)，也无限次为 \(0\)，不收敛。二进端点
至多可数，不影响“几乎处处不收敛”的结论。
\end{solution}

\begin{solution}{measure-not-lp-spike}
在 \((0,1)\) 上取
\[
 f_n=n^{2/p}\mathbf1_{(0,1/n)}.
\]
对固定 \(\varepsilon>0\)，超水平集最终为 \((0,1/n)\)，测度趋零，
故依测度收敛于零；但
\[
 \|f_n\|_p^p=n^2\cdot\frac1n=n\to\infty.
\]
\end{solution}

\begin{solution}{ae-not-measure-infinite}
对每个固定 \(x\)，当 \(n>x\) 时 \(x\notin[n,\infty)\)，故逐点趋零。
但对 \(0<\varepsilon<1\)，
\[
 m(\{|\mathbf1_{[n,\infty)}|>\varepsilon\})=\infty.
\]
若 \(E\) 有限测度，则 \(E\cap[n,\infty)\downarrow\varnothing\)，由递减
连续性 \(m(E\cap[n,\infty))\to0\)，所以局部依测度收敛。
\end{solution}

\begin{solution}{lp-not-ae-subsequence-warning}
typewriter 整列没有几乎处处极限，但可按层级抽取每层第一个区间：
\(f_{2^m}=\mathbf1_{[0,2^{-m})}\)，该子列几乎处处趋零。第
\ref{chap:egorov-subsequence} 章的子列原理只断言“存在某个子列”，并未
断言原序列或每个子列都几乎处处收敛。
\end{solution}

\begin{solution}{finite-lq-lp-estimate}
记 \(M=\mu(X)\)、\(A=\|h\|_q\)。若 \(M=0\) 或 \(A=0\)，结论成立。
以下设二者均为正，并令
\[
 r=\frac qp,\qquad s=\frac q{q-p},\qquad \frac1r+\frac1s=1.
\]
对任意 \(\lambda>0\)，Young 不等式给出
\[
 |h|^p=(\lambda|h|^p)\lambda^{-1}
 \le \frac{\lambda^r|h|^q}{r}+\frac{\lambda^{-s}}s.
\]
积分后得到
\[
 \int|h|^p\le\frac{\lambda^rA^q}{r}+\frac{M\lambda^{-s}}s.
\]
取 \(\lambda=(M/A^q)^{1/(r+s)}\)，两项中去掉系数 \(1/r,1/s\) 后的量
相等，且 \(1/r+1/s=1\)。于是
\[
 \int|h|^p\le A^pM^{1-p/q}.
\]
取 \(p\) 次方根即得所需估计。任意非零常值函数都取等号，所以常数
\(M^{1/p-1/q}\) 不能改小。
\end{solution}

\begin{solution}{distribution-functions-convergence}
给定 \(\delta>0\)，有包含关系
\[
 \{|f|>a+\delta\}\setminus\{|f_n-f|>\delta\}
 \subset\{|f_n|>a\}
\subset
 \{|f|>a-\delta\}\cup\{|f_n-f|>\delta\}.
\]
取测度和 \(n\to\infty\)，得到相应的下极限与上极限估计。再令
\(\delta\downarrow0\)。集合 \(\{|f|>a+\delta\}\) 从下逼近
\(\{|f|>a\}\)，而 \(\{|f|>a-\delta\}\) 从上逼近
\(\{|f|\ge a\}\)。有限测度保证后一列可用从上连续性；再由
\(\mu(|f|=a)=0\)，两个极限都等于 \(\mu(|f|>a)\)，故中间量收敛。
\end{solution}

\begin{solution}{counterexample-table-convergence}
可取以下代表：
\[
\begin{array}{c|c}
\text{不存在的推论}&\text{反例}\\\hline
\text{依测度}\Rightarrow\text{a.e.}&\text{typewriter 列}\\
L^p\Rightarrow\text{a.e.}&\text{同上}\\
\text{a.e.}\Rightarrow L^p&
 n^{1/p}\mathbf1_{(0,1/n)}\\
\text{依测度}\Rightarrow L^p&
 n^{2/p}\mathbf1_{(0,1/n)}\\
\text{a.e.}\Rightarrow\text{依测度（无限测度）}&
\mathbf1_{[n,\infty)}\\
\text{一致}\Rightarrow L^p\text{（无限测度）}&1/n\text{ on }\mathbb R
\end{array}
\]
各行分别按超水平集测度、逐点值和 \(p\) 次积分直接计算。有限测度条件恰好
排除最后两行所利用的无限总体积。
\end{solution}
